\section{Discussion}
\label{sec:discussion}
\vspace{-0.5em}
\paragraph{Why does frontier width determine the size of the noise gain?}
On Aircraft, 100 near-identical variants sustain a wide frontier: clean
samples keep declining throughout training, corrupted labels stagnate
($s_i^{\text{rel}} \to 0$), and velocity cleanly separates the two ---
yielding a large, persistent gain.
On CUB-200 and Dogs, classes are more separable: clean samples are learned
quickly, the frontier shrinks early, and noisy samples stagnate alongside
hard clean samples, so EWL's discrimination weakens.
EWL still matches or exceeds the SFT baseline --- earlier in training on
CUB-200 ($+1.9$--$2.8\%$) and at parity on Dogs --- but the advantage does
not persist to the final epoch.
Appendix~\ref{app:mechanism} confirms this via (EMA, velocity) cluster
separation and the clean/noisy weight-separation score.
\vspace{-0.75em}
\paragraph{When to use EWL.}
EWL is appropriate when all three hold: (i)~high inter-class similarity
sustaining a wide frontier; (ii)~suspected label noise; (iii)~the $5.9\times$
step-time overhead is acceptable.
Avoid it when classes are visually separable (narrow frontier, signal
non-discriminative), under class imbalance (the velocity mechanism is
frequency-blind), or in single-pass/streaming settings (the EMA never
converges).
Cost-benefit notes, runtime diagnostics, and extended limitations are in
Appendix~\ref{app:discussion}.
